\section{Relação com trabalhos anteriores}
\label{sec:resultado-literatura}

Esta seção compara os resultados com as ferramentas e os estudos do Capítulo~\ref{cap:referencial}. A comparação tem um limite de origem: aqui há opinião e preferência declarada, enquanto parte dos trabalhos citados mediu desempenho.

A preferência dos estudantes pela versão em português e a concordância do painel com a afirmação de que palavras em português facilitam a compreensão, com média 4,38, vão na direção da escolha feita pelo Portugol Studio, que traduz as palavras-chave para que o iniciante não precise lidar com dois idiomas ao mesmo tempo \cite{univali2026}. Os dados não alcançam uma diferença de desenho entre as duas propostas. No Portugol o vocabulário em português é fixo; na plataforma ele é escolhido por quem usa. O formulário dos estudantes comparou uma versão em inglês com uma em português definida pelos autores e, por isso, não informa o efeito de o próprio aluno escolher as palavras.

\citeonline{stefik2013} compararam o desempenho de novatos em linguagens consolidadas com o obtido em uma linguagem de controle de palavras-chave aleatórias e encontraram efeito da forma sintática sobre o desempenho. Os dados desta avaliação são de outro tipo, porque não houve tarefa nem contagem de erros. Eles não confirmam nem contrariam aquele resultado, e a comparação depende do estudo de desempenho previsto no Capítulo~\ref{cap:consideracoes}.

Cinco dos oito especialistas conheciam o Portugol Studio ou o Portugol Webstudio, e mesmo assim a afirmação de que a plataforma oferece algo que as ferramentas conhecidas não oferecem teve média 4,62. O formulário não pergunta qual seria esse diferencial. A execução passo a passo, item de maior valor didático no painel, com média 4,75, também existe no Portugol Studio, que tem depurador e inspetor de variáveis \cite{univali2026}. O que este núcleo acrescenta é a exibição dos \textit{tokens} e do código intermediário, e essas fases foram pouco visitadas: dois especialistas declararam ter visto a tabela de \textit{tokens} e quatro, o código intermediário. A avaliação diz pouco, portanto, sobre o recurso que mais distingue o núcleo das ferramentas existentes.

A correção automática pelos casos de teste, com média 4,50 entre os que lecionam, segue o ciclo de submissão e veredito de juízes \textit{online} como o Beecrowd \cite{beecrowd2026faq}. A diferença é que os casos são executados com a configuração de linguagem do aluno, como descrito na Seção~\ref{sec:integracao-plataforma}. Por fim, a preocupação com a passagem para ferramentas de mercado, que dividiu o painel no item de migração, também orienta o Portugol Studio, cuja tela reproduz a organização de uma IDE profissional para que a familiaridade seja transferível \cite{univali2026}. Lá a preocupação recai sobre a interface; aqui o painel a levantou sobre a sintaxe.

\section{Limites de interpretação e continuidade}
\label{subsec:limites-interpretacao}

Além do limite já citado sobre o público, duas ressalvas afetam a leitura dos resultados. A primeira é sobre a redação em português dos itens do SUS, para a qual não há registro de comparação com uma tradução validada. A segunda é sobre a análise: escores altos dados por especialistas não comprovam aprendizagem nem redução de carga cognitiva.

Investigar a hipótese pedagógica é a etapa seguinte e exige outro desenho de estudo: comparar, na mesma plataforma, o vocabulário padrão e o vocabulário personalizado em tarefas equivalentes de entrada e saída, seleção e repetição. A ordem das condições deve ser alternada entre os participantes, e os enunciados, os casos de teste e os critérios de classificação dos erros precisam ser definidos antes da coleta, separando erros de escrita de falhas de lógica. Com isso seria possível descrever os acertos, o tempo até a primeira execução correta e a frequência de erros em cada condição. Ainda assim, uma diferença nessas medidas não comprova, sozinha, que houve aprendizagem duradoura. Esse estudo deve ter aprovação do comitê de ética da instituição antes de qualquer coleta.
\section{Síntese e confronto com os objetivos}
\label{sec:resultado-sintese}

Esta seção retoma os objetivos específicos da introdução e diz, para cada um, o que a avaliação sustenta e o que continua sem verificação.

O primeiro objetivo era especificar a gramática da linguagem e os elementos passíveis de redefinição. O painel considerou o conjunto de elementos configuráveis adequado ao propósito da ferramenta, com média 4,57 e nenhuma nota abaixo de 4. Dois dos três professores que responderam consideraram o subconjunto suficiente para uma disciplina introdutória, e o terceiro ficou no ponto neutro. As construções apontadas como ausentes não convergiram: registros tiveram duas marcações, e as demais, uma cada.

O segundo e o terceiro objetivos eram projetar e implementar as análises léxica, sintática e semântica com suporte à configuração do vocabulário. O painel avaliou bem o resultado visível dessas fases: a resposta rápida e estável teve média 4,88, a validação prévia da configuração 4,67 e a clareza das mensagens de erro 4,71. Um respondente, que deu nota máxima às mensagens, pediu que elas sugerissem a correção, recurso que o núcleo não tem. A verificação do autor mostrou que a mensagem de \textit{token} inesperado nomeia a classe do \textit{token} pela palavra canônica, mesmo quando o usuário escreveu outra. A adaptação das mensagens ao vocabulário em uso é parcial, e esse ponto não foi captado pelo instrumento.

O quarto objetivo era integrar o núcleo à interface. A afirmação de que o código intermediário exibido corresponde ao que o programa executa teve média 4,83 entre seis respondentes, sem nenhuma nota abaixo de 4, e dois marcaram ``não sei''. Só quatro dos seis que deram nota declararam ter visto o código intermediário. O resultado não contradiz a premissa de que a interface mostra o que a máquina executa, mas é uma percepção com base pequena. A sustentação dessa premissa continua sendo a construção do sistema e os testes da Seção~\ref{sec:metodologia-testes}.

O quinto objetivo era disponibilizar o núcleo em ambiente \textit{web} e verificá-lo por testes automatizados. Cinco respondentes declararam ter escrito e executado programas na plataforma \textit{web}, e nenhum relatou erro de tradução ou de execução nas respostas abertas. A verificação por testes está descrita na Seção~\ref{sec:metodologia-testes} e não depende da avaliação.

O sexto objetivo era definir o protocolo de avaliação. O protocolo foi definido e aplicado, com oito respondentes, e a ele se somou a sessão com dezoito estudantes. A aplicação mostrou limites do próprio instrumento: dois itens mudaram de bloco e passaram a ser respondidos por no máximo quatro pessoas, o termo de consentimento ficou sem item de marcação e vários recursos receberam mais notas do que usos declarados.

Quanto à hipótese pedagógica, vale a distinção da Seção~\ref{subsec:limites-interpretacao}. O que foi coletado é opinião. A maioria do painel considera plausível que a personalização reduza a barreira sintática, com média 3,88, e um respondente discorda das três premissas. Onze de dezoito estudantes acharam a versão em português mais fácil de entender, depois de uma demonstração e de um período de uso, sem tarefa medida, e o único que nunca tinha programado achou as duas igualmente difíceis. Nenhum desses dados prova a redução da barreira nem aprendizagem. Do lado desfavorável, o painel se dividiu sobre o risco de migração para uma linguagem de mercado, três de sete respondentes concordam que a personalização pode desviar a atenção do raciocínio algorítmico e a intenção de adotar a plataforma em uma disciplina introdutória ficou em 3,00 entre os que lecionam. Esses pontos têm o mesmo peso dos resultados favoráveis e orientam os encaminhamentos do Capítulo~\ref{cap:consideracoes}.
